\documentclass[10pt,a4paper]{article}

% ---------- lingua, codifica, font (testo + matematica) ----------
\usepackage[T1]{fontenc}
\usepackage[utf8]{inputenc}
\usepackage[italian]{babel}
\usepackage{mathpazo}                 % Palatino per testo e matematica
\usepackage{microtype}

% ---------- matematica ----------
\usepackage{amsmath,amssymb,amsfonts,mathtools,bm,cancel,mathrsfs,nicefrac}
\DeclareMathOperator{\sen}{sen}       % seno all'italiana
\DeclareMathOperator{\tg}{tg}
\DeclareMathOperator{\ctg}{cotg}
\DeclareMathOperator{\arcsen}{arcsen}
\DeclareMathOperator{\arctg}{arctg}
\newcommand{\dd}{\mathrm{d}}                       % differenziale
\newcommand{\vv}[1]{\vec{#1}}                      % vettore con freccia
\newcommand{\norma}[1]{\left\lvert#1\right\rvert} % modulo di vettore
\newcommand{\pedice}[1]{_{\text{#1}}}              % pedice testuale: E\pedice{in}

% ---------- unità di misura e numeri (virgola decimale) ----------
\usepackage{siunitx}
\sisetup{output-decimal-marker={,},per-mode=symbol,
         inter-unit-product=\ensuremath{{}\cdot{}},group-digits=false}

% ---------- grafici e figure ----------
\usepackage{tikz}
\usetikzlibrary{arrows.meta,calc,positioning,patterns,decorations.pathmorphing}
\usepackage{pgfplots}
\pgfplotsset{compat=1.16}
\usepackage{graphicx}
% Grafici collocati nel punto pertinente dello svolgimento.
\newenvironment{grafico}{\begin{figure}[h!]\centering}{\end{figure}}

% ---------- impaginazione ----------
\usepackage[a4paper,left=1.8cm,right=1.8cm,top=2.3cm,bottom=2.2cm]{geometry}
\usepackage{fancyhdr,lastpage,needspace,enumitem,booktabs,array,colortbl,tabularx}
\setlength{\parindent}{0pt}
\setlength{\parskip}{3pt}
\setlength{\headheight}{14pt}
\setlist[enumerate]{leftmargin=1.7em,itemsep=2pt,topsep=2pt,parsep=0pt}
\setlist[itemize]{leftmargin=1.3em,itemsep=2pt,topsep=2pt,parsep=0pt}

% ---------- colori degli indicatori ----------
\usepackage[table,dvipsnames]{xcolor}
\definecolor{anac}{HTML}{C62828}\definecolor{anabg}{HTML}{FDECEA} % rosso
\definecolor{svic}{HTML}{1A56DB}\definecolor{svibg}{HTML}{E8F0FE} % blu
\definecolor{argc}{HTML}{A67C00}\definecolor{argbg}{HTML}{FFF6CC} % giallo
\definecolor{comc}{HTML}{2E7D32}\definecolor{combg}{HTML}{E6F4EA} % verde
% Solo griglie d'esame (su richiesta): IND arancio, DAT teal
\definecolor{indc}{HTML}{E65100}\definecolor{indbg}{HTML}{FFF0E0}
\definecolor{datc}{HTML}{00796B}\definecolor{datbg}{HTML}{E0F2F1}

% Riferimenti agli indicatori: usare queste macro in tutto il documento.
\newcommand{\coloreANA}{anac}\newcommand{\nomeANA}{Analizzare}
\newcommand{\coloreSVI}{svic}\newcommand{\nomeSVI}{Sviluppare e calcolare}
\newcommand{\coloreARG}{argc}\newcommand{\nomeARG}{Argomentare}
\newcommand{\coloreCOM}{comc}\newcommand{\nomeCOM}{Comunicare}
\newcommand{\coloreIND}{indc}\newcommand{\nomeIND}{Individuare}
\newcommand{\coloreDAT}{datc}\newcommand{\nomeDAT}{Interpretare i dati}
% Griglie d'esame: matematica \renewcommand{\nomeANA}{Comprendere}, \nomeSVI{Sviluppare il processo risolutivo};
% fisica/mista: \nomeSVI{Sviluppare il processo risolutivo}, \nomeDAT{Interpretare criticamente i dati} (mista: ...rappresentare, elaborare).
\newcommand{\sigla}[1]{\textcolor{\csname colore#1\endcsname}{\textbf{#1}}}
\newcommand{\indicatore}[1]{\textcolor{\csname colore#1\endcsname}{\textbf{#1: \csname nome#1\endcsname}}}

% ---------- box colorati e chip ----------
\usepackage{tcolorbox}
\tcbuselibrary{skins,breakable}

% Box passaggio chiave: \begin{boxana}[Titolo] ... \end{boxana}  (ana|svi|arg|com)
\newcommand{\defbox}[4]{%
  \newtcolorbox{box#1}[1][Passaggio chiave]{enhanced,breakable,frame hidden,
    borderline west={2.5pt}{0pt}{#2},colback=#3,boxrule=0pt,arc=1.5pt,
    left=7pt,right=5pt,top=2pt,bottom=2pt,before skip=4pt,after skip=4pt,
    fontupper=\small,before upper={\textbf{\idevidenza}\enspace\sigla{#4}\enspace\textbf{##1}\par}}}
\defbox{ana}{anac}{anabg}{ANA}
\defbox{svi}{svic}{svibg}{SVI}
\defbox{arg}{argc}{argbg}{ARG}
\defbox{com}{comc}{combg}{COM}
\defbox{ind}{indc}{indbg}{IND}
\defbox{dat}{datc}{datbg}{DAT}

% Box enunciato
\newtcolorbox{enun}{colback=black!3,colframe=black!25,boxrule=0.4pt,arc=2pt,
  left=6pt,right=6pt,top=3pt,bottom=3pt,fontupper=\small,
  before skip=5pt,after skip=5pt}

% Chip identificativo e massimo in punti; i dati sono forniti dal registro approvato.
\newcommand{\chip}[4]{\tcbox[on line,boxsep=0pt,left=3pt,right=3pt,top=1pt,bottom=1pt,
  arc=4pt,boxrule=0.4pt,colframe=#1,colback=#2,fontupper=\scriptsize\bfseries\color{#1}]{#3\ #4}}
\newcommand{\ANA}[2]{\chip{anac}{anabg}{\sigla{ANA}}{#1\, / massimo #2\,p}}
\newcommand{\SVI}[2]{\chip{svic}{svibg}{\sigla{SVI}}{#1\, / massimo #2\,p}}
\newcommand{\ARG}[2]{\chip{argc}{argbg}{\sigla{ARG}}{#1\, / massimo #2\,p}}
\newcommand{\COM}[2]{\chip{comc}{combg}{\sigla{COM}}{#1\, / massimo #2\,p}}
\newcommand{\IND}[2]{\chip{indc}{indbg}{\sigla{IND}}{#1\, / massimo #2\,p}}
\newcommand{\DAT}[2]{\chip{datc}{datbg}{\sigla{DAT}}{#1\, / massimo #2\,p}}
\newcommand{\ess}{\ensuremath{\bigstar}}
\newcommand{\legenda}{\indicatore{ANA}\quad\indicatore{SVI}\par
  \indicatore{ARG}\quad\indicatore{COM}}
% Griglie d'esame: \legendaesame{ANA}{IND}{SVI}{ARG} (matematica) o \legendaesame{ANA}{SVI}{DAT}{ARG} (fisica, mista)
\newcommand{\legendaesame}[4]{\indicatore{#1}\quad\indicatore{#2}\par\indicatore{#3}\quad\indicatore{#4}}
% Limite categorico sull'intera verifica; il registro non incrementa il conteggio.
\newcounter{numeroevidenze}
\newcommand{\controllaevidenze}{\stepcounter{numeroevidenze}%
  \ifnum\value{numeroevidenze}>40\relax
    \PackageError{correttore}{Superato il limite categorico di 40 evidenze}{Accorpare le evidenze prima di generare il correttore.}%
  \fi}
% Metadati mantenuti nel sorgente; nel corpo il titolo viene emesso solo dal box.
\newcommand{\idevidenza}{}
\newcommand{\evidenza}[5]{\controllaevidenze\def\idevidenza{#1}\needspace{8\baselineskip}}
% Registro raggruppato: esercizio, argomento comune, righe ID/nome/punto/indicatore/massimo.
\newcommand{\registrogruppo}[3]{\needspace{6\baselineskip}\par\smallskip{\small
  \textbf{Esercizio #1}\enspace #2\par
  \begin{tabularx}{\linewidth}{@{}lXllr@{}}\toprule
  ID & Evidenza & Punto & Ind. & Max\\\midrule
  #3\bottomrule\end{tabularx}\par}}
\newcommand{\inizioesercizi}{\clearpage}
\newcommand{\riepiloghi}[2]{\sez{Riepiloghi}
  \noindent\begin{minipage}[t]{0.58\linewidth}#1\end{minipage}\hfill
  \begin{minipage}[t]{0.39\linewidth}#2\end{minipage}\par}

\newcommand{\livelli}[5]{\begin{description}[leftmargin=2.7cm,style=nextline,font=\small\bfseries]
\item[Non raggiunto] #1
\if\relax\detokenize{#2}\relax\else\item[Iniziale] #2\fi
\item[Base] #3
\if\relax\detokenize{#4}\relax\else\item[Intermedio] #4\fi
\item[Avanzato] #5
\end{description}}

% Intestazioni di colonna colorate
\newcommand{\hANA}{\cellcolor{anabg}\sigla{ANA}}
\newcommand{\hSVI}{\cellcolor{svibg}\sigla{SVI}}
\newcommand{\hARG}{\cellcolor{argbg}\sigla{ARG}}
\newcommand{\hCOM}{\cellcolor{combg}\sigla{COM}}
\newcommand{\hIND}{\cellcolor{indbg}\sigla{IND}}
\newcommand{\hDAT}{\cellcolor{datbg}\sigla{DAT}}

% ---------- macro di struttura ----------
% Sezione semplice
\newcommand{\sez}[1]{\needspace{4\baselineskip}\par\vspace{1em}%
  {\large\bfseries #1}\par\vspace{-0.4em}\rule{\linewidth}{0.6pt}\par\vspace{0.2em}}
% Esercizio: \esercizio{Titolo}{ANA}{SVI}{ARG}{COM}{TOT}
\newcommand{\esercizio}[6]{\needspace{8\baselineskip}\par\vspace{1em}%
  {\large\bfseries #1}\par\vspace{0.25em}}
% Punto: \punto{a) titolo}
\newcommand{\punto}[1]{\needspace{6\baselineskip}\par\vspace{0.7em}{\bfseries #1}\par\vspace{0.2em}}
% Tabella punteggi del punto: \tabpunti{ANA}{SVI}{ARG}{COM}{TOT}  (usare 0 se assente)
\newcommand{\tabpunti}[5]{\par\noindent\renewcommand{\arraystretch}{1.15}\small
  \begin{tabular}{*{5}{|>{\centering\arraybackslash}p{1.4cm}}|}\hline
  \hANA & \hSVI & \hARG & \hCOM & \textbf{Tot}\\\hline
  #1 & #2 & #3 & #4 & \textbf{#5}\\\hline\end{tabular}\par\smallskip}
% Griglie d'esame: \tabpuntiesame{ANA}{IND}{SVI}{ARG}{v1}{v2}{v3}{v4}{TOT} (sigle nell'ordine dei valori)
\newcommand{\tabpuntiesame}[9]{\par\noindent\renewcommand{\arraystretch}{1.15}\small
  \begin{tabular}{*{5}{|>{\centering\arraybackslash}p{1.4cm}}|}\hline
  \csname h#1\endcsname & \csname h#2\endcsname & \csname h#3\endcsname & \csname h#4\endcsname & \textbf{Tot}\\\hline
  #5 & #6 & #7 & #8 & \textbf{#9}\\\hline\end{tabular}\par\smallskip}
% Nota (alternative, errori trascinati, criteri)
\newcommand{\nota}[1]{\begin{itemize}[font=\footnotesize]\item {\footnotesize #1}\end{itemize}}
\newcommand{\criteri}[1]{\par\nopagebreak[3]\smallskip\needspace{3\baselineskip}%
  {\footnotesize\textbf{Alternative, errori e criteri}\par
  \begin{itemize}[itemsep=1pt,topsep=1pt]#1\end{itemize}}}

% ---------- intestazione e piè di pagina ----------
\newcommand{\materia}{}
\newcommand{\classe}{}
\newcommand{\argomenti}{}
\pagestyle{fancy}\fancyhf{}
\renewcommand{\headrulewidth}{0.3pt}
\lhead{\small\textsc{Correttore\ \materia}}
\rhead{}
\cfoot{\small Pag.\ \thepage\ di \pageref{LastPage}}

\renewcommand{\materia}{Fisica}
\renewcommand{\classe}{Classe quarta, liceo scientifico \textperiodcentered\ Prova modellata sulla seconda prova di fisica}
\renewcommand{\nomeSVI}{Sviluppare il processo risolutivo}\renewcommand{\nomeDAT}{Interpretare criticamente i dati}
\begin{document}
\begin{center}
{\LARGE\bfseries Correttore\ \materia}\\[3pt]
\classe
\end{center}

\legendaesame{ANA}{SVI}{DAT}{ARG}
\sez{Registro delle evidenze}
\registrogruppo{1}{MOL}{
E01 & Riconosce la legge di Hooke come modello & intero & \sigla{ANA} & 60\\
E02 & Rappresenta i dati con barre di errore & intero & \sigla{DAT} & 80\\
E03 & Esegue la regressione e ricava $k$ con incertezza & intero & \sigla{DAT} & 100\\
E04 & Riporta $k$ con unità e cifre corrette & intero & \sigla{SVI} & 60\\
E05 & Valuta la pertinenza del modello & intero & \sigla{ARG} & 100\\
}
\registrogruppo{2}{OSC}{
E06 & Riconosce il moto armonico & intero & \sigla{ANA} & 90\\
E07 & Calcola pulsazione e periodo & intero & \sigla{SVI} & 120\\
E08 & Calcola velocità massima ed energia & intero & \sigla{SVI} & 120\\
E09 & Confronta energia e lavoro & intero & \sigla{DAT} & 70\\
E10 & Giustifica la coerenza dei risultati & intero & \sigla{ARG} & 100\\
E11 & Imposta la legge oraria & intero & \sigla{ANA} & 100\\
}
\riepiloghi{%
\begin{tabular}{@{}lrrrrr@{}}\toprule
 & \hANA & \hSVI & \hDAT & \hARG & Tot\\\midrule
Esercizio 1 & 60 & 60 & 180 & 100 & 400\\
Esercizio 2 & 190 & 240 & 70 & 100 & 600\\
\midrule
Totale & 250 & 300 & 250 & 200 & 1000\\\bottomrule
\end{tabular}}{%
\begin{tabular}{@{}lr@{}}\toprule
Argomento & Max\\\midrule
MOL, Legge di Hooke e dati sperimentali & 400\\
OSC, Oscillatore armonico & 600\\
\midrule
Totale & 1000\\\bottomrule
\end{tabular}}

\inizioesercizi
\esercizio{Esercizio 1 -- Problema, parte sperimentale}{}{}{}{}{}
\begin{enun}
Una molla è allungata da forze $F=0{,}5;\,1{,}0;\,1{,}5;\,2{,}0;\,2{,}5\ \si{N}$ (errore $\SI{0.01}{N}$) e si misurano gli allungamenti $x=21;\,40;\,62;\,81;\,101\ \si{mm}$ (errore $\SI{1}{mm}$). (a) Rappresenta i dati $F$--$x$ con le barre di errore. (b) Determina la costante elastica con una regressione lineare e la sua incertezza. (c) Discuti se la legge di Hooke è pertinente ai dati.
\end{enun}
\tabpuntiesame{ANA}{SVI}{DAT}{ARG}{60}{60}{180}{100}{400}
\evidenza{E01}{1}{ANA}{MOL}{60}
\begin{boxana}[Riconosce la legge di Hooke come modello]
\begin{enumerate}
\item La forza elastica è proporzionale all'allungamento: $F=kx$.
\item I dati vanno confrontati con una retta per l'origine.
\end{enumerate}
\end{boxana}
\evidenza{E02}{1}{DAT}{MOL}{80}
\begin{boxdat}[Rappresenta i dati con barre di errore]
\begin{enumerate}
\item Grafico $F$--$x$ con assi etichettati, unità e scala.
\item Barre di errore $\pm\SI{0.01}{N}$ e $\pm\SI{1}{mm}$.
\end{enumerate}
\end{boxdat}
\evidenza{E03}{1}{DAT}{MOL}{100}
\begin{boxdat}[Esegue la regressione e ricava $k$ con incertezza]
\begin{enumerate}
\item Regressione lineare $F=kx+b$ con il software.
\item $k=(24{,}9\pm0{,}3)\ \si{N/m}$; $b$ compatibile con zero.
\end{enumerate}
\end{boxdat}
\evidenza{E04}{1}{SVI}{MOL}{60}
\begin{boxsvi}[Riporta $k$ con unità e cifre corrette]
\begin{enumerate}
\item $k$ ha l'unità $\si{N/m}$ e due cifre dopo la virgola dell'incertezza coerenti.
\end{enumerate}
\end{boxsvi}
\evidenza{E05}{1}{ARG}{MOL}{100}
\begin{boxarg}[Valuta la pertinenza del modello]
\begin{enumerate}
\item L'intercetta è compatibile con zero e i punti stanno entro le barre di errore.
\item Il modello è pertinente nell'intervallo misurato; oltre il limite elastico potrebbe non valere.
\end{enumerate}
\end{boxarg}
\esercizio{Esercizio 2 -- Problema, parte teorica}{}{}{}{}{}
\begin{enun}
Alla molla è appesa una massa di $\SI{0.500}{kg}$ che oscilla verticalmente con ampiezza $\SI{5.0}{cm}$; trascura gli attriti e usa la costante trovata. Scrivi la legge oraria, determina periodo, velocità massima ed energia meccanica, e verifica la coerenza fra l'energia e il lavoro necessario per allungare la molla di $\SI{5.0}{cm}$.
\end{enun}
\tabpuntiesame{ANA}{SVI}{DAT}{ARG}{190}{240}{70}{100}{600}
\evidenza{E06}{2}{ANA}{OSC}{90}
\begin{boxana}[Riconosce il moto armonico]
\begin{enumerate}
\item La forza di richiamo $F=-kx$ proporzionale all'allungamento dà un moto armonico semplice.
\end{enumerate}
\end{boxana}
\evidenza{E07}{2}{SVI}{OSC}{120}
\begin{boxsvi}[Calcola pulsazione e periodo]
\begin{enumerate}
\item $\omega=\sqrt{k/m}=\sqrt{\dfrac{\SI{24.9}{N/m}}{\SI{0.500}{kg}}}=\SI{7.06}{rad/s}$.
\item $T=2\pi/\omega=\SI{0.890}{s}$.
\end{enumerate}
\end{boxsvi}
\evidenza{E08}{2}{SVI}{OSC}{120}
\begin{boxsvi}[Calcola velocità massima ed energia]
\begin{enumerate}
\item $v_{\max}=A\omega=\SI{0.353}{m/s}$.
\item $E=\tfrac12kA^2=\SI{0.0311}{J}$.
\end{enumerate}
\end{boxsvi}
\evidenza{E09}{2}{DAT}{OSC}{70}
\begin{boxdat}[Confronta energia e lavoro]
\begin{enumerate}
\item Il lavoro per allungare la molla di $A$ è $\tfrac12kA^2=\SI{0.0311}{J}$, uguale all'energia $E$.
\end{enumerate}
\end{boxdat}
\evidenza{E10}{2}{ARG}{OSC}{100}
\begin{boxarg}[Giustifica la coerenza dei risultati]
\begin{enumerate}
\item In assenza di attriti l'energia meccanica si conserva: l'energia cinetica massima $\tfrac12mv_{\max}^2$ è uguale a $E$.
\item Il valore di $T$ è compatibile con quello atteso per $k$ e $m$ di questo ordine di grandezza.
\end{enumerate}
\end{boxarg}
\evidenza{E11}{2}{ANA}{OSC}{100}
\begin{boxana}[Imposta la legge oraria]
\begin{enumerate}
\item $x(t)=A\cos(\omega t+\varphi)$ con $A=\SI{5.0}{cm}$ e fase iniziale scelta dalle condizioni iniziali.
\end{enumerate}
\end{boxana}
\end{document}
